(1) Newton径向方程为 \(\ddot r=-\mu/r^2\)。每单位质量能量 \(\mathcal E=\dot r^2/2-\mu/r\) 守恒；释放点 \(\dot r(0)=0\) 使 \(\mathcal E=-\mu/r_0\)。因为只有中心力、角动量零，运动保持严格径向。

(2) 守恒式给 \(\dot r^2=2\mu(1/r-1/r_0)\)。下落分支须取负根 \(\boxed{\dot r=-\sqrt{2\mu(1/r-1/r_0)}}\)，速率随半径缩小而增加。只凭平方根正值不能确定运动方向。

(3) 用 \(r=r_0\cos^2\eta\)，\(0\le\eta<\pi/2\)，有 \(dr=-2r_0\sin\eta\cos\eta\,d\eta\)，而速率为 \(\sqrt{2\mu/r_0}\tan\eta\)。积分得 \(\boxed{t(\eta)=\sqrt{r_0^3/(2\mu)}[\eta+\sin\eta\cos\eta]}\)。令 \(\eta\to\pi/2\)，理论坠中心时间为 \(\boxed{t_c=\frac\pi2\sqrt{r_0^3/(2\mu)}}\)。

(4) 探测器动能为 \(K=m\mu(1/r-1/r_0)\)，对时间微分 \(\dot K=-m\mu\dot r/r^2=F_r\dot r>0\)，与重力做功一致。\(r\to0\) 时牛顿点源模型预测速度无限大，但真实天体有有限半径和内部质量分布，且相对论修正可能变得重要，不能用奇点极限当真实撞击速度。
